\subsection{II.2 Problems and Failure Modes}

\paragraph{1. Failure-blind and unbounded routing.}
The outer controller always invokes the adaptive reviewer, even after the final
or a failed step, but \texttt{route\_after\_review} then terminates on failure
before consulting \texttt{needs\_adaptation}. It therefore cannot
decompose, replace or reassign precisely the step for which feedback is most
valuable. In the other direction, a successful step may trigger a replacement
tail of arbitrary length. \texttt{maximum\_number\_of\_steps\_in\_plan} appears
only in initial-planning prompts and is neither a schema constraint nor an outer
execution budget. Repeated extensions can prevent completion and produce
unbounded cost. This contradicts the study guide's central control principle:
feedback loops need an orchestration-enforced attempt budget, rather than relying
on an LLM to stop itself.

\paragraph{2. A lossy, uncalibrated feedback signal.}
The adaptive reviewer observes only the latest text summary, a small outcome
dictionary and the remaining steps. It does not receive the full completed
history, execution verdicts, workspace manifest or artefact contents that are
available elsewhere. Consequently its closed-loop observation is a lossy proxy
for the environment: a summary may omit an invalid assumption or corrupt file,
causing a false negative, while a stylistic surprise may cause an unnecessary
rewrite. A single stochastic critic supplies both the binary trigger and its
justification, with no confidence, deterministic invariant or independent check.
This is weaker than the study guide's two-gate distinction between ``did it run?''
and ``did it achieve the goal?'', and makes control vulnerable to shared model
blind spots.

\paragraph{3. Unverified and irreversible revisions.}
The initial plan receives fixed-round critique, whereas an adaptive tail is
generated once and spliced into execution without review. Pydantic guarantees
field types and permitted agent names, but not task coverage, dependency order,
artefact availability or a maximum length. Moreover, the completed prefix is
always immutable. If new evidence invalidates an earlier result, the controller
can only build on or work around it, not roll back the affected step and its
dependants. Finally, \texttt{run\_step} checkpoints before the subsequent merge;
the revised plan and \texttt{adaptive\_history} are not persisted by the merge.
A crash can therefore restore the pre-revision plan, undermining the audit trail
and reproducibility expected from an observable closed-loop system.
